\chapter{Symmetry Diagonalises}\label{linear:ch:symmetry}

\begin{center}\itshape
Which eigenvectors can we find before the matrix entries are known?\\
Start with a matrix whose rows are cyclic shifts of one row.
\end{center}

Take a row of numbers, move its last entry to the front, and use the shifted row as the next row of a matrix. Repeat until you have a square matrix. The result is \emph{circulant}. Because every row is a shifted copy of the first, this matrix has an exact symmetry.

An eigenbasis of a circulant matrix can be written down before we know its entries: use the Fourier vectors made from roots of unity. Changing the entries changes the eigenvalues, but not this eigenvector basis. That is what it means for symmetry to diagonalise a whole family of matrices at once.

The same idea reappears in Cayley graphs and in the Hamming cube, where the right basis turns distances into simple sums. Once we have those sums, linear algebra yields bounds for codes and explains the factorisation of the group determinant.

Keep the main example close: whenever you see a matrix built from translates of one pattern, try a Fourier vector before looking for a characteristic polynomial.

\section{Circulant matrices}

We work over $\C$, with the Hermitian product
$\langle x,y\rangle=\sum_jx_j\overline{y_j}$, linear in its first variable,
as in Chapter~11. Write $M^*=\overline M^{\mathsf T}$ for the adjoint.
A matrix $U$ is \emph{unitary} when $U^*U=UU^*=I$; its columns are then
an orthonormal basis. A matrix $M$ is \emph{normal} when $MM^*=M^*M$. A \emph{matrix algebra}
here means a complex vector subspace containing $I$ and closed under matrix
multiplication.

Let $n\ge1$ and $\omega=e^{2\pi i/n}$. Index vectors of $\C^n$ by $\Z/n$, and let $S$ be the cyclic shift, $(Sx)_j=x_{j+1}$, so
that $S_{j,j+1}=1$ (indices modulo $n$) and $S^n=I$. A \emph{circulant} matrix is a matrix whose rows are successive cyclic shifts
of the first one:
\[
C=\begin{pmatrix}c_0&c_1&\cdots&c_{n-1}\\c_{n-1}&c_0&\cdots&c_{n-2}\\\vdots&&\ddots&\vdots\\c_1&c_2&\cdots&c_0\end{pmatrix},
\qquad C_{jk}=c_{k-j},
\]
that is, $C=\sum_{k=0}^{n-1}c_kS^k=p(S)$ with $p(z)=c_0+c_1z+\dots+c_{n-1}z^{n-1}$, the \emph{representer polynomial}.

\begin{theorem}[eigenvalues of circulants]\label{linear:thm:circulant}
For $a\in\Z/n$ let $v_a=(\omega^{aj})_{j\in\Z/n}$. Then $Cv_a=p(\omega^a)v_a$ for every circulant $C=p(S)$. The vectors
$v_a/\sqrt n$ form an orthonormal basis of $\C^n$, so every circulant is normal, and
\[
\det C=\prod_{a=0}^{n-1}p(\omega^a).
\]
\end{theorem}
\begin{proof}
We begin with the simplest circulant: the cyclic shift $S$. Fix $a\in\{0,\ldots,n-1\}$ and write
\[
v_a=(1,\omega^a,\omega^{2a},\ldots,\omega^{(n-1)a})^{\mathsf T}.
\]
Shifting the coordinates once gives
\[
(Sv_a)_j=(v_a)_{j+1}=\omega^{a(j+1)}
=\omega^a(v_a)_j
\]
with indices taken modulo $n$; the last coordinate is consistent because $\omega^{an}=1$. Hence $Sv_a=\omega^av_a$. Repeated application gives $S^kv_a=\omega^{ak}v_a$.

Every circulant is a polynomial $C=p(S)$ in that shift, so
\[
Cv_a=\left(\sum_{k=0}^{n-1}c_kS^k\right)v_a
=\left(\sum_{k=0}^{n-1}c_k\omega^{ak}\right)v_a
=p(\omega^a)v_a.
\]
Thus the eigenvalue depends only on the representer polynomial evaluated at an $n$-th root of unity.

To show that these vectors form a basis, calculate their inner products using the standard Hermitian product on $\C^n$. For $a\ne b$ modulo $n$,
\[
\langle v_a,v_b\rangle
=\sum_{j=0}^{n-1}\omega^{(a-b)j}
=\frac{1-(\omega^{a-b})^n}{1-\omega^{a-b}}=0.
\]
For $a=b$, every entry has modulus $1$, so $\|v_a\|^2=n$. Therefore the $n$ vectors $v_a/\sqrt n$ form an orthonormal basis.

Let $U$ be the unitary matrix with these vectors as columns. The eigenvector identities say
\[
U^*CU=\operatorname{diag}\bigl(p(1),p(\omega),\ldots,p(\omega^{n-1})\bigr).
\]
If $C=UDU^*$ with $D$ diagonal, then $CC^*=UDD^*U^*=UD^*DU^*=C^*C$,
so $C$ is normal. Its determinant is the product of its eigenvalues:
\[
\det C=\prod_{a=0}^{n-1}p(\omega^a).
\]
This also explains why \emph{the same} Fourier basis diagonalises every circulant: the basis belongs to the shift $S$, not to the particular coefficients $c_k$.
\end{proof}

\paragraph{Why this works.}
A circulant matrix does not care which position we call the first. If $S$ shifts coordinates once, then a circulant $C$ is a polynomial in $S$; it therefore preserves every eigenspace of $S$. Roots of unity provide the eigenvectors of that shift. For example,
\[
C=\begin{pmatrix}2&1&1\\1&2&1\\1&1&2\end{pmatrix}
=I+J
\]
has eigenvalue $4$ on the constant vector $(1,1,1)$, and eigenvalue $1$ on the plane where coordinates sum to zero. The Fourier vectors choose a distinguished basis inside that plane. The general circulant theorem is this observation with all $n$ roots of unity written out.

The matrix $F=\frac1{\sqrt n}(\omega^{jk})_{j,k}$ is the \emph{Fourier matrix}, and the theorem says that \emph{the discrete Fourier
transform diagonalises every circulant at once}. Coordinates in this
eigenbasis are given by $F^*x$. Figure~\ref{linear:fig:circulant} shows the eigenvalues as the values of $p$ along the
unit circle. Three consequences are immediate: circulants commute with each other; the inverse of an invertible circulant is a
circulant (it is $F\operatorname{diag}(1/p(\omega^a))F^*$, a polynomial in $S$ by Lagrange interpolation); and a circulant is
invertible exactly when $p$ has no $n$-th root of unity as a root.
The coefficient row of a product is the cyclic convolution
$e_j=\sum_{r+s\equiv j\bmod n}c_r d_s$. To compute it in $O(n\log n)$
complex arithmetic operations, choose a power of two $m\ge2n-1$, with $m<4n$,
and pad the two coefficient lists by zeros to length $m$. The radix-two
algorithm proved in Chapter~11 computes their Fourier transforms in
$O(m\log m)$ operations. Multiply corresponding values and use Fourier
inversion. Since the ordinary product has degree at most $2n-2<m$, the result
has no cyclic wrap at length $m$; combine coefficients whose indices differ
by $n$ to recover $e_j$. The case $n=1$ is one multiplication. This is an
arithmetic-operation count, without a claim about rounding or bit complexity.

\begin{figure}[ht]
\centering
\fig{linear-ch09-circulant}
\caption{The circulant with first row $(1,2,0,1,0,0,0)$ has representer $p(z)=1+2z+z^3$. Its seven eigenvalues are the values
of $p$ at the seventh roots of unity: seven points on the closed curve $\theta\mapsto p(e^{i\theta})$.}
\label{linear:fig:circulant}
\end{figure}

\begin{example}[the $3\times3$ circulant]\label{linear:exa:circ3}
With $\omega=e^{2\pi i/3}$, the determinant of the circulant with first row $(a,b,c)$ is
\[
(a+b+c)(a+b\omega+c\omega^2)(a+b\omega^2+c\omega)=a^3+b^3+c^3-3abc.
\]
The second and third factors are complex conjugate when $a,b,c$ are real, with product $a^2+b^2+c^2-ab-bc-ca=\frac12\bigl((a-b)^2+
(b-c)^2+(c-a)^2\bigr)$. This factor is positive unless $a=b=c$.
Thus for real $a,b,c$ not all equal, the cubic has the sign of $a+b+c$,
including value zero when that sum is zero. When $a=b=c$, the cubic is zero
regardless of the sum. Exercise~\ref{linear:ex:circ3} finds every integer
of this form.
\end{example}

\begin{example}[cycles and their relatives]\label{linear:exa:cycle}
For $n\ge3$, the simple cycle $C_n$ has adjacency matrix $S+S^{-1}$.
Its eigenvalues follow here from the Fourier vectors:
$\omega^a+\omega^{-a}=2\cos(2\pi a/n)$. For $n=1,2$, that matrix instead
counts a loop twice or an edge twice; it is not a simple-cycle adjacency
matrix. Any simple graph on $\Z/n$ in which $j\sim k$ depends only on $k-j$, a \emph{circulant graph},
has the eigenvalues $\sum_{s\in D}\omega^{as}$, where $D$ is the set of allowed differences. For $K_n$, $D=\Z/n\setminus\{0\}$ and
the eigenvalues are $n-1$ and, for $a\ne0$, $\sum_{s\ne0}\omega^{as}=-1$.
\end{example}

Conversely, symmetry forces the shape:

\begin{proposition}\label{linear:prop:commutant}
A matrix commutes with $S$ if and only if it is a circulant.
\end{proposition}
\begin{proof}
Index rows and columns by $\mathbb Z/n\mathbb Z$.
By definition, a circulant matrix has the form
$M_{jk}=c_{k-j}$ for a sequence $c_0,\ldots,c_{n-1}$.
Equivalently $M$ is a polynomial in the cyclic
shift $S$, so every circulant commutes with $S$.

For the converse suppose $MS=SM$.
Our shift convention is $(Sf)(j)=f(j+1)$,
with indices taken modulo $n$.
A direct multiplication of entries gives
\[
(MS)_{jk}=M_{j,k-1},
\qquad (SM)_{jk}=M_{j+1,k}.
\]
Equality of the two products therefore means
$M_{j,k-1}=M_{j+1,k}$. Replace $k$ by $k+1$
to obtain
\[
M_{jk}=M_{j+1,k+1}.
\]
Repeating this identity around the cycle
shows that the entry depends only on the
difference $k-j\pmod n$. Indeed,
$M_{jk}=M_{0,k-j}$ after $n-j$ simultaneous
shifts. Set $c_t=M_{0t}$; then
$M_{jk}=c_{k-j}$, which is the circulant form.
Thus the whole commutant of $S$ consists
precisely of circulant matrices.
\end{proof}

\section{Characters of finite abelian groups}

Circulants are the case $G=\Z/n$ of a construction that works for every finite abelian group $G$, written additively. Let $\C^G$
be the space of functions $G\to\C$ with the inner product $\langle f,h\rangle=\sum_{x\in G}f(x)\overline{h(x)}$. For $g\in G$, the
\emph{translation} $T_g$ acts by $(T_gf)(x)=f(x+g)$. The translations commute, $T_gT_h=T_{g+h}$, and each one is unitary (it
permutes the values).

\begin{definition}\label{linear:def:character}
A \emph{character} of $G$ is a group homomorphism $\chi:G\to\C^\times$, that is, $\chi(x+y)=\chi(x)\chi(y)$. The characters form
the \emph{dual group} $\widehat G$ under pointwise multiplication. The trivial character is $\mathbf 1$.
\end{definition}

If $N=|G|$, then $\chi(x)^N=\chi(Nx)=\chi(0)=1$, so the values of a character are $N$-th roots of unity, $|\chi(x)|=1$ and
$\chi^{-1}=\bar\chi$. For $G=\Z/n$, a character is determined by $\chi(1)$, an $n$-th root of unity, so the characters are the
$\chi_a(x)=\omega^{ax}$: the vectors $v_a$ of Theorem~\ref{linear:thm:circulant}.

Chapter~11 proved that there are $|G|$ characters and that they form an
orthogonal basis. Here we will construct that basis by simultaneously
diagonalising the translations, without using the abelian structure theorem.

\begin{lemma}\label{linear:lem:commuting}
A family of pairwise commuting diagonalisable operators on a finite-dimensional complex space has a common basis of eigenvectors.
\end{lemma}
\begin{proof}
Induct on $d=\dim_{\mathbb C}V$. The case $d=0$
is empty and the case $d=1$ is immediate.
If every operator in the family is scalar,
any basis is a common eigenbasis.
Otherwise choose a non-scalar diagonalizable
operator $T$ from the family. It has at least
two distinct eigenvalues, so
\[
V=\bigoplus_{\lambda\in\operatorname{spec}T}
E_\lambda(T)
\]
is a direct sum of non-zero proper subspaces.

For every other operator $U$ in the family,
commutativity with $T$ implies
$T(Uv)=U(Tv)=\lambda Uv$ when $v\in E_\lambda(T)$.
Thus each eigenspace of $T$ is $U$-invariant.
The restricted operators still commute.
They are also diagonalizable: the minimal
polynomial of $U|_{E_\lambda(T)}$ divides
that of $U$, and the latter is a product
of distinct linear factors over $\mathbb C$.

Each summand has dimension less than $d$,
so the induction hypothesis gives a common
eigenbasis for the restricted family on
each $E_\lambda(T)$. Concatenating these
bases gives a basis for $V$, and every
vector in it is an eigenvector of $T$ and
of every $U$. This also handles infinite
commuting families: no enumeration or finite
generating subset of the family was required.
\end{proof}

\begin{theorem}[characters of a finite abelian group]\label{linear:thm:characters}
Let $G$ be a finite abelian group.
\begin{enumerate}[label=(\roman*)]
\item For distinct characters $\chi\ne\psi$, $\langle\chi,\psi\rangle=0$; and $\langle\chi,\chi\rangle=|G|$.
\item There are exactly $|G|$ characters, and they form an orthogonal basis of $\C^G$.
\item The eigenvectors common to all translations are exactly the non-zero multiples of characters: $T_g\chi=\chi(g)\chi$.
\end{enumerate}
\end{theorem}
\begin{proof}
\emph{(i) Orthogonality.}
A character takes its values among roots of unity,
because $\chi(g)^{|G|}=\chi(|G|g)=1$.
In particular $|\chi(g)|=1$. If $\chi\ne\psi$,
then $\varphi(g)=\chi(g)\overline{\psi(g)}$
is a non-trivial character. Choose $h$ with
$\varphi(h)\ne1$. Translating the index
of the finite sum gives
\[
\sum_{g\in G}\varphi(g)
=\sum_{g\in G}\varphi(g+h)
=\varphi(h)\sum_{g\in G}\varphi(g),
\]
so the sum is zero. This is
$\langle\chi,\psi\rangle=0$.
When $\chi=\psi$ every summand equals one,
hence $\langle\chi,\chi\rangle=|G|$.

\emph{(iii) Common eigenvectors.}
For a character,
\[
(T_g\chi)(x)=\chi(x+g)=\chi(g)\chi(x).
\]
Conversely, suppose $f\ne0$ satisfies
$T_gf=c(g)f$ for every $g\in G$.
Evaluating at $x=0$ yields
$f(g)=c(g)f(0)$. If $f(0)=0$,
this would give $f=0$, so $f(0)\ne0$.
Scale $f$ to have $f(0)=1$; then $f(g)=c(g)$.
Because $T_{g+h}=T_gT_h$,
\[
c(g+h)f=T_{g+h}f=c(g)c(h)f,
\]
and $c(0)=1$. Thus $c$ is a character
and $f$ is a non-zero multiple of it.

\emph{(ii) Completeness.}
Each translation $T_g$ has finite order
dividing $|G|$; its minimal polynomial
divides $X^{|G|}-1$, whose roots over
$\mathbb C$ are all distinct.
Thus every $T_g$ is diagonalizable.
The translations commute, so
Lemma~\ref{linear:lem:commuting} supplies
a common eigenbasis of the $|G|$-dimensional
space $\mathbb C^G$. By (iii), each basis
vector is a multiple of a character,
and these characters are distinct because
the vectors are independent. We have
at least $|G|$ characters. By (i),
distinct characters are non-zero and
mutually orthogonal, so there can be
at most $|G|$ in $\mathbb C^G$.
Hence there are exactly $|G|$, and they
form the desired orthogonal basis.
\end{proof}

Arrange the normalized character columns in the square matrix
$U=(\chi(x))_{x,\chi}/\sqrt{|G|}$. Orthogonality gives $U^*U=I$; a square
matrix with a left inverse is invertible, so $UU^*=I$ too. Its row identities
give \emph{dual orthogonality}:
$\sum_{\chi\in\widehat G}\chi(x)\overline{\chi(y)}=|G|\,[x=y]$.
This appendix uses the \emph{unnormalized Fourier transform}
$\hat f(\chi)=\langle f,\chi\rangle=\sum_xf(x)\overline{\chi(x)}$;
Chapter~11 put $1/|G|$ in the forward transform instead. Here expansion in the
orthogonal basis is the \emph{Fourier inversion formula}
\[
f=\frac1{|G|}\sum_{\chi\in\widehat G}\hat f(\chi)\,\chi.
\]

Now the guiding principle. Call a matrix $M=(M_{x,y})_{x,y\in G}$ \emph{translation-invariant} if $M_{x+g,y+g}=M_{x,y}$ for all
$x,y,g$, or equivalently $M_{x,y}=m(x-y)$ for a function $m$ on $G$; equivalently again, $M$ commutes with every $T_g$.
Indeed, simultaneous translation by $-y$ gives $M_{x,y}=M_{x-y,0}$.
Also $(MT_g)_{x,y}=M_{x,y-g}$ and $(T_gM)_{x,y}=M_{x+g,y}$, so equality
of these entries is exactly simultaneous translation invariance.

\begin{theorem}[symmetry diagonalises]\label{linear:thm:invariant}
Every translation-invariant matrix $M_{x,y}=m(x-y)$ on a finite abelian group $G$ has every character as an eigenvector:
\[
M\chi=\hat m(\chi)\,\chi,\qquad\hat m(\chi)=\sum_{u\in G}m(u)\overline{\chi(u)}.
\]
In particular $\det M=\prod_{\chi\in\widehat G}\hat m(\chi)$.
\end{theorem}
\begin{proof}
Choose a character $\chi$ of $G$.
By the translation-invariance hypothesis,
\[
(M\chi)(x)=\sum_{y\in G}m(x-y)\chi(y).
\]
Make the bijective substitution $u=x-y$.
Since $\chi$ is multiplicative on the
additive group,
$\chi(x-u)=\chi(x)\chi(-u)$.
Every character value has modulus one,
so $\chi(-u)=\chi(u)^{-1}
=\overline{\chi(u)}$. Therefore
\[
\begin{aligned}
(M\chi)(x)
 &=\sum_{u\in G}m(u)\chi(x-u)\\
 &=\chi(x)\sum_{u\in G}m(u)
          \overline{\chi(u)}
 =\widehat m(\chi)\chi(x).
\end{aligned}
\]
Thus $\chi$ is an eigenvector with
the claimed eigenvalue. By
Theorem~\ref{linear:thm:characters}, the
$|G|$ characters form a basis of
$\mathbb C^G$. In this common basis,
the matrix of $M$ is diagonal,
with diagonal entries
$\widehat m(\chi)$ as $\chi$
ranges over the character group.
Its determinant is the product of
those entries, proving
\[
\det M=\prod_{\chi\in\widehat G}
\widehat m(\chi).
\]
The character basis is independent
of the particular function $m$:
that is why all translation-invariant
matrices can be diagonalised together.
\end{proof}

For $G=\Z/n$ this is Theorem~\ref{linear:thm:circulant} again (with $m(u)=c_{-u}$). The chosen eigenbasis does not
depend on $m$: all translation-invariant matrices are diagonalised by the same basis, and their eigenvalues are Fourier
coefficients. For two such matrices, put $v=x-z$ in their product entry:
\[
(MM')_{x,y}=\sum_zm(x-z)m'(z-y)
=\sum_vm(v)m'(x-y-v).
\]
Thus their kernel is the \emph{convolution} $(m*m')(u)=\sum_vm(v)m'(u-v)$.
On each character the eigenvalues multiply, giving
$\widehat{m*m'}=\hat m\,\hat m'$. Exercise~\ref{linear:ex:convolution}
also proves the corresponding norm identity.

\begin{remark}
For a non-abelian group, written multiplicatively, put $L_g e_x=e_{gx}$ and
$R_h e_x=e_{xh^{-1}}$. Left and right translations commute, but if $ab\ne ba$,
then $L_aL_b\ne L_bL_a$, as their values at $e_e$ show. Thus the matrices
commuting with all right translations cannot all be diagonal in one basis.
The regular-representation decomposition proved in Chapter~11 gives blocks
of dimension $d$, each repeated $d$ times, for every irreducible of degree $d$.
At the end of this appendix we use it to prove the corresponding determinant
factorization, including irreducibility of its factors.
\end{remark}

\section{Cayley graphs and the Paley graphs}

Let $D\subseteq G\setminus\{0\}$ be \emph{symmetric}, $D=-D$. The \emph{Cayley graph} $\operatorname{Cay}(G,D)$ has vertex set $G$,
with $x\sim y$ when $y-x\in D$. It is $|D|$-regular, and its adjacency matrix is translation-invariant with $m=\mathbf 1_D$. By
Theorem~\ref{linear:thm:invariant}, its eigenvalues are
\[
\lambda_\chi=\sum_{s\in D}\overline{\chi(s)}=\sum_{s\in D}\chi(s),
\]
real because $D=-D$ pairs $\chi(s)$ with $\chi(-s)=\overline{\chi(s)}$. Circulant graphs are the case $G=\Z/n$. The hypercube
$Q_n$ is $\operatorname{Cay}((\Z/2)^n,\{e_1,\dots,e_n\})$, and its characters are $\chi_T(x)=(-1)^{\sum_{i\in T}x_i}$; we come back to
it in the next section.

The next example uses squares. Let $p$ be a prime with $p\equiv1\pmod4$, and let $D\subseteq\Z/p$ be the set of non-zero
squares. It has $\frac{p-1}2$ elements (each non-zero square is $x^2$ for exactly two values $\pm x$), and it is symmetric because
$-1$ is a square modulo $p$ when $p\equiv1\pmod4$ (Volume~2, Chapter~6). The \emph{Paley graph} $P_p$ is $\operatorname{Cay}(\Z/p,D)$:
two residues are friends when their difference is a square. Figure~\ref{linear:fig:paley} shows $P_{13}$.
A simple graph is \emph{strongly regular with parameters} $(v,k,\lambda,\mu)$
if it has $v$ vertices, each has $k$ neighbours, and two distinct vertices
have $\lambda$ common neighbours when adjacent and $\mu$ when not adjacent.

\begin{figure}[ht]
\centering
\fig{linear-ch09-paley}
\caption{The Paley graph $P_{13}$: $x\sim y$ when $x-y$ is one of the squares $\pm1,\pm3,\pm4$ modulo $13$ (difference $\pm1$ charcoal,
$\pm3$ rust, $\pm4$ ochre). It is strongly regular with parameters $(13,6,2,3)$ and eigenvalues $6$ and $\frac{-1\pm\sqrt{13}}2$.}
\label{linear:fig:paley}
\end{figure}

\begin{theorem}[Paley graphs]\label{linear:thm:paley}
Let $p\equiv1\pmod4$ be prime and $\zeta=e^{2\pi i/p}$. The eigenvalues of $P_p$ are $\frac{p-1}2$ (once) and $\frac{-1\pm\sqrt p}2$,
each with multiplicity $\frac{p-1}2$. Consequently $P_p$ is strongly regular with parameters $\bigl(p,\frac{p-1}2,\frac{p-5}4,
\frac{p-1}4\bigr)$.
\end{theorem}
\begin{proof}
Let $D$ be the set of non-zero quadratic residues
modulo $p$, so the graph is the Cayley graph
$\operatorname{Cay}(\mathbb Z/p,D)$.
The additive characters $\chi_a(x)=\zeta^{ax}$
diagonalise its adjacency matrix $A$, with
eigenvalues
\[
\lambda_a=\sum_{s\in D}\zeta^{as}.
\]
For $a=0$ this is $|D|=(p-1)/2$.

For $a\ne0$ consider the Gauss sum
$g(a)=\sum_{x\in\mathbb F_p}\zeta^{ax^2}$.
Each non-zero square has two square roots,
while zero has one; hence
$g(a)=1+2\lambda_a$.
To find its absolute value, use the
bijective change of variables $x=y+t$:
\[
\begin{aligned}
|g(a)|^2
 &=\sum_{x,y\in\mathbb F_p}
          \zeta^{a(x^2-y^2)}\\
 &=\sum_{t\in\mathbb F_p}\zeta^{at^2}
       \sum_{y\in\mathbb F_p}\zeta^{2aty}=p.
\end{aligned}
\]
Indeed, when $t=0$ the inner sum is $p$,
and for $t\ne0$ it is the vanishing sum
of a non-trivial additive character,
since $2at\ne0$. As $p\equiv1\pmod4$,
$-1$ is a quadratic residue. Choose
$c$ with $c^2=-1$. Replacing $x$ by
$cx$ proves
\[
\overline{g(a)}
=\sum_x\zeta^{-ax^2}
=\sum_x\zeta^{a(cx)^2}=g(a).
\]
The Gauss sum is real, and its square
is $p$, so $g(a)=\pm\sqrt p$. Therefore
\[
\lambda_a=\frac{-1\pm\sqrt p}{2}
\qquad(a\ne0).
\]
The diagonal of $A$ is zero, hence
$\operatorname{tr}A=0$. Summing the
eigenvalues and substituting
$\lambda_a=(g(a)-1)/2$ gives
$\sum_{a\ne0}g(a)=0$.
All these sums equal $\sqrt p$
or $-\sqrt p$; therefore each sign
occurs exactly $(p-1)/2$ times.

It remains to read the common-neighbour
counts from the spectrum. Set
$k=(p-1)/2$ and
$\theta,\tau=(-1\pm\sqrt p)/2$.
They satisfy $\theta+\tau=-1$ and
$\theta\tau=-(p-1)/4$. Since $A$ is
real symmetric, its character eigenbasis
is orthogonal, and
\[
B=(A-\theta I)(A-\tau I)
  =A^2+A-\frac{p-1}{4}I
\]
vanishes on $\mathbf1^\perp$.
Thus $B=cJ$ for some scalar $c$.
Evaluating on $\mathbf1$, where
$A\mathbf1=k\mathbf1$, gives
\[
cp=k^2+k-\frac{p-1}{4}
=\frac{p(p-1)}4,
\qquad c=\frac{p-1}4.
\]
Consequently
\[
A^2=\frac{p-1}4J-A+\frac{p-1}4I.
\]
For distinct vertices, $(A^2)_{xy}$
counts common neighbours: it equals
$(p-5)/4$ when $x\sim y$, and
$(p-1)/4$ otherwise. On the diagonal
it equals $k$, as required. These are
exactly the asserted strongly regular
parameters.
\end{proof}

For $p=5$, $P_5$ is the pentagon, with eigenvalues $2$ and $\frac{-1\pm\sqrt5}2$: the golden ratio, once more. For $p=13$ we get
$\mathrm{srg}(13,6,2,3)$. Which sign goes with which $a$? Exercise~\ref{linear:ex:gausssign} shows that $g(a)=\bigl(\frac ap\bigr)g(1)$,
so the answer depends only on the sign of $g(1)$. Gauss conjectured in 1801 that $g(1)=+\sqrt p$, and it took him four more years
to prove it; in his letter of 3 September 1805 to Olbers, he compared the solution to a flash of lightning.

We can earn the needed spectral bound here. Let a real symmetric adjacency
matrix on $v$ vertices have constant degree $k$ and least eigenvalue $\tau<0$.
For an independent set of size $s>0$, its indicator is
$x=(s/v)\mathbf1+y$ with $y\perp\mathbf1$ and $\|y\|^2=s-s^2/v$.
The real spectral theorem of Volume~1, Chapter~7 gives
$y^{\mathsf T}Ay\ge\tau\|y\|^2$. No two selected vertices are adjacent, so
\[
0=x^{\mathsf T}Ax\ge k\frac{s^2}{v}
  +\tau\left(s-\frac{s^2}{v}\right).
\]
Dividing by $s>0$ and $k-\tau>0$ yields
$s\le-v\tau/(k-\tau)$; an empty set satisfies it too.
This is \emph{Hoffman's ratio bound}. Write $\alpha(P_p)$ for the largest
independent-set size. The eigenvalues just proved give
\[
\alpha(P_p)\le p\,\frac{-\tau}{k-\tau}=p\,\frac{(1+\sqrt p)/2}{(p+\sqrt p)/2}=\sqrt p.
\]
\emph{A set of residues in which no difference is a square has at most $\sqrt p$ elements.} Multiplying by a non-square maps
$P_p$ onto its complement: multiplication swaps the two cosets of squares
and nonsquares among nonzero residues, and hence swaps edges and nonedges.
Thus the same bound holds for sets in which every nonzero difference is a
square. Exercise~\ref{linear:ex:paleyself} checks the isomorphism and its
spectral consequence in detail.

\begin{remark}
\emph{Outlook.} The bound proved here is not the end of the subject.
Hanson and Petridis's article in the 2021 journal volume obtains the stronger prime-field
clique bound $\omega(P_p)(\omega(P_p)-1)\le(p-1)/2$ by an auxiliary-polynomial
method. Their introduction records earlier small improvements for certain
primes, so this was not the first improvement of any kind. That stronger
research theorem is further reading; no argument here uses it.
\end{remark}

\section{The Hamming cube and Krawtchouk polynomials}

Now let $n\ge0$ and $G=(\Z/2)^n$, the words of length $n$ over $\{0,1\}$ with addition modulo $2$, and write $|x|$ for the \emph{weight} of $x$
(its number of ones) and $d(x,y)=|x+y|$ for the \emph{Hamming distance}. The characters are the $2^n$ functions
\[
\chi_T(x)=(-1)^{\sum_{i\in T}x_i},\qquad T\subseteq\{1,\dots,n\}.
\]
They are characters, there are $2^n$ of them, so by Theorem~\ref{linear:thm:characters} there are no others. The matrix $H_n=(\chi_T(x))
_{x,T}$ of size $2^n$ is the \emph{Walsh--Hadamard matrix}: its entries are $\pm1$ and $H_n^{\mathsf T}H_n=2^nI$. It is the
Kronecker power $H_1^{\otimes n}$ of $H_1=\bigl(\begin{smallmatrix}1&1\\1&-1\end{smallmatrix}\bigr)$:
for $n\ge1$, in first-bit order, $\chi_{(t,T')}(b,x')=(-1)^{tb}\chi_{T'}(x')$ gives the
block matrix $\bigl(\begin{smallmatrix}H_{n-1}&H_{n-1}\\H_{n-1}&-H_{n-1}\end{smallmatrix}\bigr)$.
Here $\otimes$ means replacing each entry of the first matrix by that scalar
times the second matrix, and $H_0=[1]$. Compute the two half-transforms and
combine their outputs by sums and differences. The operation count satisfies
$T(n)=2T(n-1)+2^n$, $T(0)=0$, so induction gives $T(n)=n2^n$ additions
and subtractions. Exercise~\ref{linear:ex:walsh} writes out the recursion.

The cube has more symmetry than translations: permuting coordinates preserves distances. The matrices invariant under both are
the functions of the distance alone. For $0\le k\le n$, let $D_k$ be the \emph{distance-$k$ matrix}, $(D_k)_{x,y}=1$ if $d(x,y)=k$
and $0$ otherwise. Thus $D_0=I$ and $\sum_kD_k=J$; for $n\ge1$, $D_1$ is the adjacency matrix of $Q_n$.

\begin{definition}\label{linear:def:krawtchouk}
For a variable $t$, put $\binom tj=t(t-1)\cdots(t-j+1)/j!$, with
$\binom t0=1$. For a nonnegative integer top entry this agrees with the
usual binomial coefficients, including value zero when the lower entry
exceeds the top entry.
The \emph{Krawtchouk polynomials} for length $n$ are
\[
K_k(t)=\sum_{j=0}^k(-1)^j\binom tj\binom{n-t}{k-j},\qquad k=0,1,\dots,n.
\]
\end{definition}

So $K_0=1$, and when the indices are available,
$K_1(t)=n-2t$ and $K_2(t)=\binom{n-t}2-t(n-t)+\binom t2$.
Every summand defining $K_k$ has degree at most $k$. Its coefficient of $t^k$
is $(-1)^k/(j!(k-j)!)$, so the total leading coefficient is
$(-1)^k\sum_{j=0}^k1/(j!(k-j)!)=(-2)^k/k!\ne0$.
For integers $0\le t\le n$, expanding the two finite binomials gives
\[
\sum_{k=0}^n K_k(t)z^k=(1-z)^t(1+z)^{n-t}:
\]
the coefficient of $z^k$ chooses $j$ of the $t$ negative factors and $k-j$
of the $n-t$ positive factors. At $z=1$ this gives
$\sum_k K_k(t)=2^n[t=0]$, treating the empty factor at $t=0$ as $1$.
Exercise~\ref{linear:ex:krawgen} also extracts the reflection identity.

\begin{theorem}[distances in the cube]\label{linear:thm:krawtchouk}
\begin{enumerate}[label=(\roman*)]
\item $D_k\chi_T=K_k(|T|)\,\chi_T$ for every $T$. The spectrum is the list
$K_k(t)$, $0\le t\le n$, each repeated $\binom nt$ times; when values
coincide, their repetitions are combined to give the eigenvalue's multiplicity.
\item For every word $z$, the characters of weight $k$ sum to $\sum_{|T|=k}\chi_T(z)=K_k(|z|)$.
\item (Orthogonality) For all $k,l$,
\[
\sum_{t=0}^n\binom ntK_k(t)K_l(t)=2^n\binom nk[k=l].
\]
\end{enumerate}
\end{theorem}
\begin{proof}
\emph{(i) Distance-matrix eigenvalues.}
View the cube as the additive group
$(\mathbb Z/2\mathbb Z)^n$.
The matrix $D_k$ has entries
$1_{\{|x-y|=k\}}$ and is therefore
translation-invariant. By
Theorem~\ref{linear:thm:invariant},
the character $\chi_T(x)=(-1)^{|T\cap x|}$
has eigenvalue
\[
\sum_{|u|=k}\chi_T(u)
=\sum_{|u|=k}(-1)^{|T\cap u|}.
\]
Write $t=|T|$. A vector $u$ of weight
$k$ whose support meets $T$ in $j$
positions can be chosen in
$\binom tj\binom{n-t}{k-j}$ ways.
Its contribution is $(-1)^j$, so
the eigenvalue equals
\[
\sum_j(-1)^j\binom tj\binom{n-t}{k-j}
=K_k(t).
\]
Exactly $\binom nt$ characters have weight $t$. They account for the stated
list of eigenvalues; equal values from different weights contribute to the
same eigenspace and their multiplicities add.

\emph{(ii) Character sum by weight.}
Fix $z$ of weight $t$ instead, and
sum $\chi_T(z)=(-1)^{|T\cap z|}$
over sets $T$ of size $k$.
The same intersection count yields
$\sum_{|T|=k}\chi_T(z)=K_k(t)$.

\emph{(iii) Orthogonality.}
Calculate $\operatorname{tr}(D_kD_l)$
first in the standard basis:
\[
\operatorname{tr}(D_kD_l)
=\sum_{x,y}(D_k)_{xy}(D_l)_{yx}.
\]
The two entries are simultaneously
one precisely when the same pair
$(x,y)$ has distance both $k$ and
$l$. For $k\ne l$ the sum is zero;
for $k=l$, each of the $2^n$
vertices has $\binom nk$ neighbours
at distance $k$, so it is
$2^n\binom nk$. On the other hand,
the common character basis
diagonalises both matrices, and
the same trace is
\[
\sum_{T\subseteq[n]}K_k(|T|)K_l(|T|)
=\sum_{t=0}^n\binom ntK_k(t)K_l(t).
\]
Equating the two expressions proves
the stated orthogonality identity.
\end{proof}

The $n+1$ matrices $D_0,\dots,D_n$ span a commutative algebra of symmetric
matrices. For closure, simultaneous translations and coordinate permutations
preserve each $D_k$, hence preserve each product after reindexing its middle
summation variable. These permutations take any pair of words of distance $i$
to any other such pair: translate the first word to zero and permute the
support of the second. Thus a product's entries depend only on distance, so
it is a linear combination of the $D_i$. They commute because the character
basis diagonalises them all. This is called the \emph{Bose--Mesner algebra}
of the Hamming distance partition. The Krawtchouk values are its eigenvalue table.

\section{Delsarte's linear programming bound}

A \emph{binary code} of length $n\ge0$ is a set $C\subseteq\{0,1\}^n$.
For an integer $d\ge1$, say its minimum distance is \emph{at least} $d$
when every distinct pair has distance at least $d$; this includes empty and
one-word codes. A code with this property corrects
$\lfloor(d-1)/2\rfloor$ errors: two balls of that radius cannot meet, since
the triangle inequality would make their centres less than $d$ apart.
That triangle inequality holds coordinate by coordinate: a difference between
two endpoint bits requires a difference on at least one step through a third bit.
Write $A(n,d)$ for the largest size of a length-$n$ code with minimum distance
at least $d$. The maximum exists because the cube is finite. If $d>n$, it
equals $1$; for $d=0$, separately, all $2^n$ words are allowed.
Constructions give lower bounds; upper bounds need an argument that works for \emph{every} code, including
codes with no structure at all.

In his 1973 thesis at the Universit\'e Catholique de Louvain, published as a
Philips Research Reports supplement, Philippe Delsarte developed such an
argument from the linear algebra of the previous section. Its starting point is the \emph{distance distribution} of a code,
\[
A_i=\frac1{|C|}\#\{(x,y)\in C^2:d(x,y)=i\},\qquad i=0,\dots,n,
\]
which we use only for $C\ne\varnothing$; an empty code satisfies all size bounds
directly. It is the average number of codewords at distance $i$ from a codeword.
Then $A_0=1$, $A_i\ge0$, $\sum_iA_i=|C|$, and $A_i=0$ for
$1\le i<d$. Delsarte saw that the $A_i$ satisfy many more, hidden, inequalities.

\begin{theorem}[Delsarte, 1973]\label{linear:thm:delsarte}
For every nonempty code $C\subseteq\{0,1\}^n$, $n\ge0$, and every $k=0,1,\dots,n$,
\[
\sum_{i=0}^nA_i\,K_k(i)\ge0.
\]
\end{theorem}
\begin{proof}
Recall that the distance distribution is
\[
A_i=\frac1{|C|}
  |\{(x,y)\in C^2:d(x,y)=i\}|.
\]
In the binary cube, subtraction is addition
and $d(x,y)=|x+y|$. By
Theorem~\ref{linear:thm:krawtchouk}(ii),
\[
K_k(|x+y|)
=\sum_{|T|=k}\chi_T(x+y).
\]
Every cube character is real-valued
($\chi_T=\pm1$) and multiplicative,
so $\chi_T(x+y)=\chi_T(x)\chi_T(y)$.
Substitute these identities and
interchange finite sums:
\[
\begin{aligned}
\sum_{i=0}^nA_iK_k(i)
&=\frac1{|C|}\sum_{x,y\in C}K_k(|x+y|)\\
&=\frac1{|C|}\sum_{|T|=k}
   \sum_{x,y\in C}\chi_T(x)\chi_T(y)\\
&=\frac1{|C|}\sum_{|T|=k}
   \left(\sum_{x\in C}\chi_T(x)\right)^2.
\end{aligned}
\]
Each square is a non-negative real
number, and $|C|>0$. Therefore the
sum is non-negative for every $k$.
The inequality comes from a
sum-of-squares identity, not from
an estimate or an asymptotic argument.
\end{proof}

In matrix language: $K_k(d(x,y))$ is the matrix $E_k=\sum_{|T|=k}\chi_T\chi_T^{\mathsf T}$, a positive semidefinite matrix (a multiple of
the orthogonal projection onto the span of the characters of weight $k$), and the inequality says $\mathbf 1_C^{\mathsf T}E_k\mathbf
1_C\ge0$. Whatever the code, its indicator vector cannot have a negative ``energy'' in any eigenspace.

\begin{corollary}[the linear programming bound]\label{linear:cor:lp}
For integers $n\ge0$ and $d\ge1$, $A(n,d)$ is at most the maximum of
$1+\sum_{i=d}^na_i$ over all real $a_d,\dots,a_n\ge0$ satisfying
\[
\binom nk+\sum_{i=d}^na_iK_k(i)\ge0\qquad(k=1,\dots,n).
\]
\end{corollary}
\begin{proof}
The maximum in the statement exists. The choice $a_i=0$ is feasible.
Sum all $n$ displayed inequalities. The generating identity just proved gives
$\sum_{k=1}^nK_k(i)=-1$ for $i\ge d\ge1$, while
$\sum_{k=1}^n\binom nk=2^n-1$. Hence $\sum_i a_i\le2^n-1$.
The feasible coordinates are non-negative and bounded; any maximizing sequence
has a coordinatewise convergent subsequence by the finite compactness argument
in Appendix~\ref{linear:ch:perron}. Its limit preserves the finite inequalities
and attains the maximum. If $d>n$ there are no variables and the maximum is $1$.

Take any nonempty binary code $C$ of length $n$
and minimum distance at least $d$.
Its distance distribution satisfies
$A_0=1$, $A_i=0$ for $1\le i<d$,
and $A_i\ge0$ for every $i$.
Moreover, the Delsarte inequalities
(Theorem~\ref{linear:thm:delsarte}) state
\[
\sum_{i=0}^n A_iK_k(i)\ge0
\qquad (k=1,\ldots,n).
\]
Since $K_k(0)=\binom nk$, substituting
$a_i=A_i$ for $i\ge d$ shows that the
numbers arising from $C$ satisfy every
constraint in the stated linear program.
The sum of the distance-distribution
coefficients is
\[
1+\sum_{i=d}^n a_i
=\sum_{i=0}^nA_i=|C|.
\]
Thus the objective value of \emph{every}
code is bounded by the maximum feasible
objective value. Taking the largest
possible code size gives the claimed
upper bound on $A(n,d)$.
\end{proof}

The objective and constraints are linear in finitely many real variables:
this is a \emph{linear program}. Even one of its inequalities can give a
useful bound. Here is the case $k=1$.

\begin{corollary}[Plotkin bound]\label{linear:cor:plotkin}
For integers $n\ge0$, $d\ge1$ with $2d>n$, every code of length $n$ and
minimum distance at least $d$ has $|C|\le\frac{2d}{2d-n}$.
\end{corollary}
\begin{proof}
An empty code satisfies the bound directly. If $n=0$, there is only one word
and the bound is $1$. Otherwise let $C$ be nonempty, of minimum distance
at least $d$ and distance distribution
$A_i$. The degree-one Krawtchouk
polynomial is $K_1(i)=n-2i$.
Delsarte's inequality for $k=1$ says
\[
0\le\sum_{i=0}^n A_i(n-2i)
 =n+\sum_{i=d}^n A_i(n-2i),
\]
using $A_0=1$ and $A_i=0$ for
$0<i<d$. Since $i\ge d>n/2$,
each $n-2i$ is at most the
strictly negative number $n-2d$.
Therefore
\[
0\le n+\sum_{i=d}^n A_i(n-2i)
 \le n-(2d-n)\sum_{i=d}^n A_i.
\]
The denominator $2d-n$ is positive,
so we may rearrange without reversing
the inequality:
\[
|C|-1=\sum_{i=d}^n A_i
\le\frac{n}{2d-n},
\qquad
|C|\le\frac{2d}{2d-n}.
\]
\end{proof}

\begin{example}[a perfect fit]\label{linear:exa:simplex}
For $n=7$, $d=4$, Plotkin gives $|C|\le8$. The eight words
\[
0000000,\ 1110100,\ 0111010,\ 0011101,\ 1001110,\ 0100111,\ 1010011,\ 1101001
\]
(the zero word and the seven cyclic shifts of $1110100$) are pairwise at distance exactly $4$, so $A(7,4)=8$. Equality in the proof
forces all non-zero distances to be equal to $d$ (Exercise~\ref{linear:ex:plotkinsharp}), which is the case here.
\end{example}

\begin{remark}
The sum-of-squares identity works for every code, whether or not its words
form a subspace. It supplies an exact certificate: a non-negative combination
of the inequalities, with coefficients we can check by hand, can bound every
feasible distance distribution. The final exercise uses just the first two
inequalities to settle $A(8,4)$. The papers listed below develop asymptotic code
bounds and an analytic counterpart for sphere packings; those outlooks supply
no premise for the finite bounds proved here.
\end{remark}

\section{Group determinants}

The determinant formula for circulants extends to every finite abelian group.

\begin{theorem}[Dedekind]\label{linear:thm:groupdet}
Let $G$ be a finite abelian group and $(x_g)_{g\in G}$ independent variables. Then
\[
\det\bigl(x_{g-h}\bigr)_{g,h\in G}=\prod_{\chi\in\widehat G}\Bigl(\sum_{g\in G}\chi(g)\,x_g\Bigr).
\]
\end{theorem}
\begin{proof}
First assign arbitrary complex numbers to the
variables $x_g$. The resulting numerical
matrix $M_{g,h}=x_{g-h}$ is
translation-invariant. Applying
Theorem~\ref{linear:thm:invariant}
to the function $m(u)=x_u$ gives
\[
\det M=
\prod_{\chi\in\widehat G}
\left(\sum_{u\in G}x_u\overline{\chi(u)}\right).
\]
This identity holds for \emph{every}
complex assignment of the $x_g$.
Both sides are polynomials in the
finitely many variables $x_g$,
so they are identical in
$\mathbb C[(x_g)_{g\in G}]$.
For justification, a one-variable polynomial vanishing at every complex
number is zero, by the root bound. Induct on the number of variables:
fix all but the last, conclude that every coefficient polynomial vanishes
at every assignment of the remaining variables, and apply induction.

Finally, $\overline\chi$ is itself
a character of $G$, since complex
conjugation preserves multiplication
of roots of unity. Conjugating twice
returns $\chi$, so
$\chi\mapsto\overline\chi$ permutes
the finite character group.
Reindex the product by this bijection:
\[
\det(x_{g-h})_{g,h\in G}
=\prod_{\chi\in\widehat G}
\left(\sum_{u\in G}\chi(u)x_u\right).
\]
No complex conjugation acts on the
variables here; it merely relabels
the factors indexed by characters.
\end{proof}

For $G=\Z/3$ this is Example~\ref{linear:exa:circ3}; for $G=(\Z/2)^2$ it says that
\[
\det\begin{pmatrix}a&b&c&d\\b&a&d&c\\c&d&a&b\\d&c&b&a\end{pmatrix}=(a+b+c+d)(a+b-c-d)(a-b+c-d)(a-b-c+d).
\]
Each linear factor is the eigenvalue on one character; its coefficients are
that character's row of the character table.

\paragraph{The non-abelian factorization.}
We can also earn the block version mentioned earlier. Let $G$ now be any
finite group, written multiplicatively, and take one representative $\rho$
of each complex irreducible representation, of degree $d_\rho$.
For independent variables $x_g$, put $A_\rho(x)=\sum_g x_g\rho(g)$. Then
\[
\det(x_{gh^{-1}})_{g,h\in G}
=\prod_\rho\det A_\rho(x)^{d_\rho},
\]
and each factor $\det A_\rho(x)$ is irreducible of degree $d_\rho$ in
$\C[(x_g)_{g\in G}]$. Here are proofs of both assertions.

The space $\C[G]$ has basis $e_g$. Give it multiplication by
$e_g e_h=e_{gh}$ and extend by distributivity and complex linearity.
Associativity follows from associativity in $G$, and $e_e$ is the identity.
Each representation extends linearly to this algebra, and
\[
\Phi:\C[G]\longrightarrow\bigoplus_\rho M_{d_\rho}(\C),\qquad
\sum_g a_g e_g\longmapsto
\left(\sum_g a_g\rho(g)\right)_\rho
\]
preserves products. If $\Phi(a)=0$, then $a$ acts as zero on every
irreducible. The decomposition of the regular representation proved in
Chapter~11 (Corollary~\ref{cor:mult}) makes left multiplication by $a$ zero
on all of $\C[G]$. Applying it to $e_e$ gives $a=0$, so $\Phi$ is injective.
The same corollary proves $|G|=\sum_\rho d_\rho^2$; domain and codomain
therefore have equal dimension, and $\Phi$ is bijective.

For $a=\sum_g x_g e_g$, left multiplication has matrix
$(x_{gh^{-1}})_{g,h}$ in the group basis. Under $\Phi$, it acts on the
$\rho$ block by $B\mapsto A_\rho(x)B$. Each of the $d_\rho$ columns of $B$
is acted on separately by $A_\rho(x)$, giving determinant
$\det A_\rho(x)^{d_\rho}$. Multiplication over the blocks proves the
displayed identity for numerical assignments, and the polynomial argument
above proves it identically.

For irreducibility, first prove that the determinant of a generic $d\times d$
matrix $Y=(y_{ij})$ is irreducible. The case $d=1$ is the variable $y_{11}$.
For $d\ge2$, expand as
$\det Y=D y_{dd}+E$ in the ring $R[y_{dd}]$, where $R$ contains all the
other variables and $D$ is the leading $(d-1)\times(d-1)$ determinant.
The rings involved are UFDs by repeated application of
Theorem~\ref{thm:gaussufd}. By induction $D$ is irreducible, still so after
adjoining unused variables: a factorization of a polynomial independent of
a variable must have both factors independent of that variable, since their
degrees in it add in a domain. Also $D$ does not divide $E$. Set the first
$d-2$ diagonal entries to $1$, all other entries to $0$ except
$y_{d-1,d}=y_{d,d-1}=1$. Then $D=0$ while $E=-1$.
Since $D$ is irreducible in the UFD $R$, this proves $\gcd(D,E)=1$.
The linear polynomial $D y_{dd}+E$ is irreducible over the fraction field
of $R$ and primitive over $R$; Proposition~\ref{prop:gaussirr} therefore
makes it irreducible over $R$ too.

Finally $\Phi$ is an invertible linear change from the $x_g$ to all entries
of all blocks $A_\rho(x)$. Those entries are consequently independent
coordinates. In these coordinates each factor is a generic determinant in
its own block variables, hence irreducible by the preceding proof, even
after adjoining the other block variables. An invertible change of variables
preserves factorization and irreducibility. The factor has degree $d_\rho$,
and factors from different blocks are not scalar multiples because they
depend on disjoint nonempty sets of independent variables. This earns the
full non-abelian factorization. For $Q_8$, the degree-two factor is the
determinant on its two-dimensional irreducible from Exercise~\ref{ex:centralchar}.

\begin{beauty}
The first row of a circulant determines its eigenvalues, while the shift
determines their common eigenvectors. On the binary cube, the same separation
turns distance counts into sums of squares. The eight words of length seven
and the sixteen affine truth tables are concrete tests: both meet an upper
bound obtained from those sums. In the next appendix, linear algebra bounds
sets of objects even when no group action is given.
\end{beauty}

\section{Exercises}

\begin{exercise}[$\star$]\label{linear:ex:circalg}
Show that the $n\times n$ circulants form a commutative algebra of dimension
$n$, closed under adjoints. Show that the inverse of every invertible
circulant is circulant, and compute $(I-tS)^{-1}$ explicitly for $t^n\ne1$.
\end{exercise}

\begin{exercise}[$\star\star$]\label{linear:ex:circ3}
Show that an integer $N$ can be written $N=a^3+b^3+c^3-3abc$ with $a,b,c\in\Z$ if and only if $N\not\equiv\pm3\pmod9$.
\end{exercise}

\begin{exercise}[$\star$]\label{linear:ex:fourier}
Show that the Fourier matrix $F$ satisfies $F^2=R$, where $R$ is the reversal $x_j\mapsto x_{-j}$, and $F^4=I$. Deduce that the
eigenvalues of $F$ are among $1,-1,i,-i$.
\end{exercise}

\begin{exercise}[$\star$]\label{linear:ex:dualgroup}
Show that $\widehat{G\times H}\cong\widehat G\times\widehat H$, with $(\chi,\psi)\mapsto\bigl((x,y)\mapsto\chi(x)\psi(y)\bigr)$, and that
$\widehat{\Z/n}\cong\Z/n$.
\end{exercise}

\begin{exercise}[$\star\star$]\label{linear:ex:convolution}
Define $(f*h)(x)=\sum_yf(y)h(x-y)$. Prove the convolution theorem $\widehat{f*h}=\hat f\,\hat h$ and Plancherel's formula
$\sum_x|f(x)|^2=\frac1{|G|}\sum_\chi|\hat f(\chi)|^2$. Show that the translation-invariant matrix with symbol $m$ acts on $f$ by
$f\mapsto m*f$.
\end{exercise}

\begin{exercise}[$\star\star$]\label{linear:ex:walsh}
With $H_0=[1]$, show that for $n\ge1$ we have $H_n=H_1\otimes H_{n-1}$ (in the block sense $\bigl(\begin{smallmatrix}H_{n-1}&H_{n-1}\\H_{n-1}&-H_{n-1}\end{smallmatrix}\bigr)$
for a suitable order of the words), and deduce an algorithm computing $H_nf$ with $n2^n$ additions and subtractions.
\end{exercise}

\begin{exercise}[$\star$]\label{linear:ex:krawgen}
Show that $\sum_{k=0}^nK_k(t)z^k=(1-z)^t(1+z)^{n-t}$ for integers $0\le t\le n$. Deduce that $K_k(n-t)=(-1)^kK_k(t)$ and that
$\sum_kK_k(t)=2^n[t=0]$.
\end{exercise}

\begin{exercise}[$\star\star$]\label{linear:ex:gausssign}
Show that $g(a)=\bigl(\frac ap\bigr)g(1)$ for $a\ne0$, where $\bigl(\frac ap\bigr)$ is the Legendre symbol. Deduce that in $P_p$ the
eigenvalue $\frac{-1+\sqrt p}2$ belongs to the characters $\chi_a$ with $a$ a square if $g(1)>0$.
\end{exercise}

\begin{exercise}[$\star$]\label{linear:ex:paleyself}
Let $p$ be a prime with $p\equiv1\pmod4$ and $c$ a non-square modulo $p$. Show that $x\mapsto cx$ is an isomorphism from $P_p$ to its complement.
Deduce again that $P_p$ and its complement have the same spectrum.
\end{exercise}

\begin{exercise}[$\star\star$]\label{linear:ex:paleyramsey}
Show that the clique number of $P_{17}$ is $3$, by hand or with the help of the bound $\sqrt{17}$. Deduce that the Ramsey number
$R(4,4)$ is at least $18$. Here $R(4,4)$ is the least vertex count forcing
a monochromatic four-clique in every red/blue colouring of the complete graph.
\end{exercise}

\begin{exercise}[$\star\star$]\label{linear:ex:plotkinsharp}
Show that a code achieving the Plotkin bound $|C|=\frac{2d}{2d-n}$ (with $2d>n$) has all its non-zero distances equal to $d$, and
that each coordinate takes the value $1$ in exactly half of its codewords.
\end{exercise}

\begin{exercise}[$\star\star\star$]\label{linear:ex:lpthree}
Plotkin's bound says nothing for $n=8$, $d=4$. Use Delsarte's inequalities for $k=1$ and $k=2$ to show that $A(8,4)\le16$, and find
a code of that size.
\end{exercise}

\section{Hints}

\begin{hint}{linear:ex:circalg}
Circulants are the polynomials in $S$ of degree less than $n$. For the inverse, expand $\sum_jt^jS^j$ and use $S^n=I$.
\end{hint}
\begin{hint}{linear:ex:circ3}
Try $(a,b,c)=(k+1,k,k)$, $(k+1,k+1,k)$, $(k+1,k,k-1)$. For the converse, reduce modulo $3$: $a^3+b^3+c^3-3abc\equiv(a+b+c)^3$.
\end{hint}
\begin{hint}{linear:ex:fourier}
$(F^2)_{jk}=\frac1n\sum_l\omega^{l(j+k)}$.
\end{hint}
\begin{hint}{linear:ex:dualgroup}
A character of $G\times H$ restricts to characters of $G\times\{0\}$ and $\{0\}\times H$.
\end{hint}
\begin{hint}{linear:ex:convolution}
Expand $\widehat{f*h}(\chi)$ and substitute $u=x-y$. For Plancherel, apply the inversion formula.
\end{hint}
\begin{hint}{linear:ex:walsh}
Split $x=(x_1,x')$ and $T$ according to whether $1\in T$.
\end{hint}
\begin{hint}{linear:ex:krawgen}
Expand $(1-z)^t(1+z)^{n-t}$ with the binomial theorem twice.
\end{hint}
\begin{hint}{linear:ex:gausssign}
Substitute $x\mapsto bx$ when $a=b^2c$.
\end{hint}
\begin{hint}{linear:ex:paleyself}
$cx-cy$ is a square if and only if $x-y$ is not.
\end{hint}
\begin{hint}{linear:ex:paleyramsey}
Use the symmetry: you may assume a clique contains $0$ and $1$.
\end{hint}
\begin{hint}{linear:ex:plotkinsharp}
Look at where the inequalities in the proof of Corollary~\ref{linear:cor:plotkin} can be equalities. For the coordinates, write
$\sum_{x,y\in C}d(x,y)$ coordinate by coordinate.
\end{hint}
\begin{hint}{linear:ex:lpthree}
The unknowns are $a_4,\dots,a_8$, and $K_2(i)=-4,-2,4,14,28$ for $i=4,\dots,8$. Add a quarter of the $k=2$ inequality to the
$k=1$ inequality. For the code, think of affine functions on $(\Z/2)^3$.
\end{hint}

\section{Solutions}

\begin{solution}{linear:ex:circalg}
\idea A circulant is a polynomial in the cyclic shift $S$, where $S^n=I$. We may therefore add and multiply circulants like polynomials, reducing exponents modulo $n$. Invertibility and adjoints can be handled with the same shift.

\writeup By Proposition~\ref{linear:prop:commutant}, every circulant has a unique expression
\[
C=c_0I+c_1S+\cdots+c_{n-1}S^{n-1},
\qquad S^n=I.
\]
The shift matrices $I,S,\ldots,S^{n-1}$ are linearly independent: in a fixed row their non-zero entries lie in different columns, so a linear combination can vanish only if every coefficient is zero. Hence the space of circulants has dimension $n$.

Sums and scalar multiples clearly remain polynomials in $S$. Products do too, since $S^jS^k=S^{j+k}$ and powers can be reduced modulo $n$. Any two polynomials in the same matrix commute, so circulants form a commutative subalgebra of $M_n(\C)$.

The shift matrix is a permutation matrix, hence unitary; its adjoint is $S^*=S^{-1}=S^{n-1}$. Consequently
\[
C^*=\sum_{j=0}^{n-1}\overline{c_j}(S^*)^j
\]
is a polynomial in $S$ and therefore another circulant.

If $C$ is an invertible circulant, it commutes with $S$. Multiplying $CS=SC$ on both sides by $C^{-1}$ shows that $C^{-1}S=SC^{-1}$. Proposition~\ref{linear:prop:commutant} then says that $C^{-1}$ is circulant.

For the requested example, use the finite geometric identity:
\[
(I-tS)(I+tS+\cdots+t^{n-1}S^{n-1})
=I-t^nS^n=(1-t^n)I.
\]
If $t^n\ne1$, the scalar on the right is non-zero, so
\[
\boxed{(I-tS)^{-1}=\frac1{1-t^n}
\sum_{j=0}^{n-1}t^jS^j}.
\]
This is an explicit circulant inverse, valid also for $t=0$.
\end{solution}

\begin{solution}{linear:ex:circ3}
\idea Factor the expression as a sum times a non-negative quadratic form. Three short choices of $(a,b,c)$ then realise every integer in the permitted congruence classes. Modulo $3$, the factorisation forces an extra factor of $3$ whenever the result is divisible by $3$.

\writeup The classical identity is
\[
a^3+b^3+c^3-3abc
=(a+b+c)(a^2+b^2+c^2-ab-bc-ca).
\]
Write the second factor as
\[
Q=\frac{(a-b)^2+(b-c)^2+(c-a)^2}{2}.
\]
For $(a,b,c)=(k+1,k,k)$, the sum is $3k+1$ and $Q=1$, giving $N=3k+1$. For $(k+1,k+1,k)$, the sum is $3k+2$ and again $Q=1$, giving $N=3k+2$. Finally, for $(k+1,k,k-1)$, the sum is $3k$ and the three squared differences are $1,1,4$, so $Q=3$ and $N=9k$. These identities work for every $k\in\Z$, including negative $k$.

Thus all integers congruent to $1$ or $2$ modulo $3$, as well as all multiples of $9$, are represented. It remains to show that integers congruent to $3$ or $6$ modulo $9$ are impossible.

Reduce the cubic identity modulo $3$. Fermat's congruence gives $x^3\equiv x\pmod3$, and $3abc\equiv0$, so
\[
N\equiv a+b+c\pmod3.
\]
If $3\mid N$, then $3\mid(a+b+c)$. But the quadratic factor can also be written
\[
Q=(a+b+c)^2-3(ab+bc+ca),
\]
which is divisible by $3$ whenever $a+b+c$ is. Hence $9\mid N=(a+b+c)Q$.

We have proved the exact criterion
\[
\boxed{N=a^3+b^3+c^3-3abc\text{ for some }a,b,c\in\Z
\iff N\not\equiv3,6\pmod9}.
\]
The excluded residues are exactly $\pm3$ modulo $9$.
\end{solution}

\begin{solution}{linear:ex:fourier}
\idea The Fourier matrix has entries made of roots of unity. Multiply it by itself and use the finite geometric sum: it vanishes unless the exponent is divisible by $n$.

\writeup Number the rows and columns by residues $0,\ldots,n-1$, and write
\[
F_{jk}=\frac1{\sqrt n}\omega^{jk},
\qquad \omega=e^{2\pi i/n}.
\]
The entry in row $j$, column $k$ of $F^2$ is
\[
\begin{aligned}
(F^2)_{jk}
&=\sum_{\ell=0}^{n-1}F_{j\ell}F_{\ell k}\\
&=\frac1n\sum_{\ell=0}^{n-1}\omega^{j\ell}\omega^{\ell k}\\
&=\frac1n\sum_{\ell=0}^{n-1}\omega^{\ell(j+k)}.
\end{aligned}
\]
If $j+k\equiv0\pmod n$, every term is $1$, so the sum equals $n$. Otherwise the ratio $\omega^{j+k}$ is a non-trivial $n$-th root of unity, and the geometric sum is
\[
\frac{1-(\omega^{j+k})^n}{1-\omega^{j+k}}=0.
\]
Thus $(F^2)_{jk}$ is $1$ exactly when $k\equiv-j\pmod n$, and zero otherwise. This is the matrix $R$ that reverses indices:
\[
(Rx)_j=x_{-j},\qquad F^2=R.
\]
Reversing twice returns the original index, so $R^2=I$. Consequently
\[
F^4=(F^2)^2=R^2=I.
\]
If $v\ne0$ is an eigenvector of $F$ with eigenvalue $\lambda$, then
\[
v=F^4v=\lambda^4v.
\]
Hence $\lambda^4=1$, so every eigenvalue lies in
\[
\{1,-1,i,-i\}.
\]
The result says that these four values are \emph{possible}; it does not assert that all four occur for every $n$. The case $n=1$ already shows why that distinction matters.
\end{solution}

\begin{solution}{linear:ex:dualgroup}
\idea A character of a product group can be restricted to each factor. Conversely, two characters can be multiplied to obtain a character on the product. Those two constructions should undo one another.

\writeup Let $G,H$ be finite abelian groups, written additively, and let $\widehat G$ denote the group of complex characters of $G$ under pointwise multiplication. Given $\chi\in\widehat G$ and $\psi\in\widehat H$, define
\[
F_{\chi,\psi}:G\times H\to\C^*,
\qquad
F_{\chi,\psi}(x,y)=\chi(x)\psi(y).
\]
This is a character: for $(x,y),(x',y')\in G\times H$,
\[
\begin{aligned}
F_{\chi,\psi}(x+x',y+y')
&=\chi(x+x')\psi(y+y')\\
&=\chi(x)\chi(x')\psi(y)\psi(y')\\
&=F_{\chi,\psi}(x,y)F_{\chi,\psi}(x',y').
\end{aligned}
\]

We obtain a homomorphism
\[
\Phi:\widehat G\times\widehat H\longrightarrow\widehat{G\times H},
\quad(\chi,\psi)\mapsto F_{\chi,\psi}.
\]
It is injective: if $F_{\chi,\psi}=F_{\chi',\psi'}$, evaluate first at $(x,0)$, where every character takes the value $1$ at $0$. We get $\chi(x)=\chi'(x)$ for every $x$, so $\chi=\chi'$. Evaluation at $(0,y)$ similarly gives $\psi=\psi'$.

It is also surjective, and we can prove this without counting. Let $F$ be \emph{any} character of $G\times H$. Define
\[
\chi(x)=F(x,0),\qquad \psi(y)=F(0,y).
\]
Both are characters of their respective groups. Since $(x,y)=(x,0)+(0,y)$,
\[
F(x,y)=F(x,0)F(0,y)=\chi(x)\psi(y)=F_{\chi,\psi}(x,y).
\]
Thus every character lies in the image, so
\[
\boxed{\widehat{G\times H}\cong\widehat G\times\widehat H}.
\]

For $G=\Z/n\Z$, a character is determined by $\chi([1])$, because $\chi([k])=\chi([1])^k$. The relation $n[1]=[0]$ requires
\[
\chi([1])^n=1.
\]
Conversely, every $n$-th root of unity $\omega^m=e^{2\pi im/n}$ defines a character $\chi_m([k])=\omega^{mk}$. Multiplication of characters corresponds to addition of exponents modulo $n$, hence
\[
\boxed{\widehat{\Z/n\Z}\cong\Z/n\Z}.
\]
This identifies the cyclic group with its dual abstractly; the particular identification depends on the chosen primitive root of unity.
\end{solution}

\begin{solution}{linear:ex:convolution}
\idea A character turns addition in a finite abelian group into multiplication of complex numbers. Substituting a sum $x=y+u$ therefore splits the Fourier transform of a convolution into a product. Orthogonality gives the energy identity.

\writeup Write the finite abelian group $G$ additively and use the \emph{unnormalised} Fourier transform
\[
\hat f(\chi)=\sum_{x\in G}f(x)\overline{\chi(x)}.
\]
For convolution $(f*h)(x)=\sum_y f(y)h(x-y)$, substitute $u=x-y$:
\[
\begin{aligned}
\widehat{f*h}(\chi)
&=\sum_x\sum_y f(y)h(x-y)\overline{\chi(x)}\\
&=\sum_y\sum_u f(y)h(u)\overline{\chi(y+u)}\\
&=\left(\sum_y f(y)\overline{\chi(y)}\right)
  \left(\sum_u h(u)\overline{\chi(u)}\right)\\
&=\boxed{\hat f(\chi)\hat h(\chi)}.
\end{aligned}
\]

Fourier inversion reads
\[
f(x)=\frac1{|G|}\sum_{\chi\in\widehat G}\hat f(\chi)\chi(x).
\]
Multiply by $\overline{f(x)}$, sum over $x$ and expand both terms. Orthogonality of characters gives
\[
\sum_{x\in G}\chi(x)\overline{\psi(x)}
=\begin{cases}|G|,&\chi=\psi,\\0,&\chi\ne\psi.\end{cases}
\]
Only equal-character terms survive, producing
\[
\boxed{\sum_{x\in G}|f(x)|^2
=\frac1{|G|}\sum_{\chi\in\widehat G}|\hat f(\chi)|^2}.
\]
The factor $1/|G|$ is essential for this unnormalised convention; a normalised transform uses a different coefficient.

Finally a translation-invariant matrix has entries $M_{xy}=m(x-y)$. Its action is
\[
(Mf)(x)=\sum_{y\in G}m(x-y)f(y)=(m*f)(x).
\]
Thus convolution is exactly the matrix operation behind translation invariance, and the Fourier transform diagonalises it.
\end{solution}

\begin{solution}{linear:ex:walsh}
\idea A Walsh character factors across coordinates, so splitting a binary word by its first bit yields a block matrix. Computing the two half-transforms and then taking their sum and difference gives the fast Walsh--Hadamard algorithm.

\writeup For $n=0$, $H_0=[1]$ returns its input without an arithmetic operation.
Now let $n\ge1$, write $x=(x_1,x')\in\mathbb F_2\times\mathbb F_2^{n-1}$
and similarly $T=(t_1,T')$. The character is
\[
\chi_T(x)=(-1)^{t_1x_1}\chi_{T'}(x').
\]
Order rows and columns first by their first bit and then by the remaining $(n-1)$ bits. If either first bit equals $0$, the extra factor is $1$; if both equal $1$, it is $-1$. Therefore
\[
\boxed{H_n=
\begin{pmatrix}
H_{n-1}&H_{n-1}\\
H_{n-1}&-H_{n-1}
\end{pmatrix}
=H_1\otimes H_{n-1}},
\quad
H_1=\begin{pmatrix}1&1\\1&-1\end{pmatrix}.
\]

For a vector $f=\binom uv$ split into two halves of size $2^{n-1}$, the block formula gives
\[
H_nf=
\begin{pmatrix}H_{n-1}(u+v)\\H_{n-1}(u-v)\end{pmatrix}.
\]
This is equivalent to first computing the two smaller transforms and then adding and subtracting their outputs. At each of the $2^{n-1}$ paired positions, forming the sum and difference requires two operations: $2^n$ additions or subtractions in all.

If $T(n)$ counts these scalar operations, then
\[
T(n)=2T(n-1)+2^n,\qquad T(0)=0.
\]
Divide by $2^n$:
\[
\frac{T(n)}{2^n}=\frac{T(n-1)}{2^{n-1}}+1.
\]
Induction gives $T(n)/2^n=n$, and hence
\[
\boxed{T(n)=n2^n}.
\]
The direct matrix multiplication would require on the order of $4^n$ operations, so the recursive factorisation is exponentially faster.
\end{solution}

\begin{solution}{linear:ex:krawgen}
\idea Expand two binomial polynomials. The coefficient of $z^k$ counts a choice of $k$ positions with a sign depending on how many came from a distinguished set of size $t$. That is precisely the Krawtchouk polynomial.

\writeup Recall
\[
K_k(t)=\sum_{j=0}^k(-1)^j
\binom tj\binom{n-t}{k-j},
\]
with binomial coefficients outside their usual range interpreted as $0$. The binomial theorem gives
\[
(1-z)^t=\sum_{j=0}^t(-1)^j\binom tjz^j,
\qquad
(1+z)^{n-t}=\sum_{\ell=0}^{n-t}\binom{n-t}{\ell}z^\ell.
\]
When the two sums are multiplied, the coefficient of $z^k$ collects terms with $j+\ell=k$, so it equals $K_k(t)$. Therefore
\[
\boxed{\sum_{k=0}^nK_k(t)z^k=(1-z)^t(1+z)^{n-t}}.
\]

Replace $t$ with $n-t$. The right-hand side becomes
\[
(1-z)^{n-t}(1+z)^t,
\]
which is the original right-hand side with $z$ replaced by $-z$. Comparing coefficients of $z^k$ yields
\[
\boxed{K_k(n-t)=(-1)^kK_k(t)}.
\]

Finally set $z=1$ in the generating function. If $t>0$, the factor $(1-z)^t$ becomes zero, so $\sum_kK_k(t)=0$. If $t=0$, the expression becomes $(1+1)^n=2^n$. Thus
\[
\boxed{\sum_{k=0}^nK_k(t)=2^n[t=0]},
\]
where $[t=0]$ equals $1$ if $t=0$ and $0$ otherwise. The case $t=0$ should be treated separately rather than assigning an ambiguous value to $0^0$.
\end{solution}

\begin{solution}{linear:ex:gausssign}
\idea Scaling the variable in a quadratic Gauss sum by a non-zero square does not change its value. Thus the sum has one value on square parameters and another on non-squares. Their total is zero, so the two values are opposites.

\writeup Let $p$ be an odd prime, $\zeta=e^{2\pi i/p}$, and
\[
g(a)=\sum_{x\in\mathbb F_p}\zeta^{a x^2},
\qquad a\in\mathbb F_p^*.
\]
If $a=b^2a'$ with $b\ne0$, substitute $y=bx$. Multiplication by $b$ is a bijection of $\mathbb F_p$, so
\[
g(a)=\sum_x\zeta^{a'(bx)^2}=\sum_y\zeta^{a'y^2}=g(a').
\]
Therefore $g(a)$ depends only on the square class of $a$. It equals $g(1)$ when $a$ is a square and one constant $g(c)$ when $a$ is a non-square.

To compare these values, sum over all non-zero $a$:
\[
\begin{aligned}
\sum_{a\ne0}g(a)
&=\sum_{x\in\mathbb F_p}\sum_{a\ne0}\zeta^{ax^2}\\
&=(p-1)+\sum_{x\ne0}(-1)=0.
\end{aligned}
\]
For $x=0$, the inner sum has $p-1$ ones. For each $x\ne0$, $a\mapsto ax^2$ permutes the non-zero residues, so the inner sum is $\sum_{u=1}^{p-1}\zeta^u=-1$.

Exactly $(p-1)/2$ non-zero values of $a$ are squares, and the same number are non-squares. Hence
\[
0=\frac{p-1}2g(1)+\frac{p-1}2g(c),
\]
so $g(c)=-g(1)$. This is precisely
\[
\boxed{g(a)=\left(\frac ap\right)g(1)}.
\]

In the Paley graph $P_p$, the non-constant character $\chi_a$ has adjacency eigenvalue
\[
\theta_a=\frac{g(a)-1}{2}.
\]
For $p\equiv1\pmod4$, if $g(1)>0$, then $g(1)=\sqrt p$ and
$\theta_a=(-1+\sqrt p)/2$ exactly when $(a/p)=1$, that is, when $a$ is a
non-zero square. For non-squares the eigenvalue is $(-1-\sqrt p)/2$.
With the fixed $\zeta=e^{2\pi i/p}$ the sign is determined, not a free choice;
this conditional conclusion does not need its evaluation.
\end{solution}

\begin{solution}{linear:ex:paleyself}
\idea An edge of a Paley graph is a pair of vertices whose difference is a non-zero square. Multiplication by a non-square exchanges the square and non-square classes, so it reverses adjacency.

\writeup Let $p\equiv1\pmod4$ be prime and let $c\in\mathbb F_p^*$ be a non-square. The subgroup $S$ of non-zero squares has index $2$ in $\mathbb F_p^*$, and its other coset is $cS$, the set of non-squares.

The Paley graph $P_p$ has vertex set $\mathbb F_p$, and two distinct vertices $x,y$ are adjacent precisely when $x-y\in S$. Define
\[
\varphi:\mathbb F_p\to\mathbb F_p,\qquad\varphi(x)=cx.
\]
Since $c\ne0$, this map is a bijection. Moreover
\[
\varphi(x)-\varphi(y)=c(x-y).
\]
If $x-y$ is a non-zero square, its product with the non-square $c$ is a non-square; if $x-y$ is a non-square, its product with $c$ is a square. Consequently
\[
\{x,y\}\in E(P_p)
\iff
\{\varphi(x),\varphi(y)\}\notin E(P_p)
\]
for every $x\ne y$. Thus $\varphi$ is an isomorphism
\[
\boxed{P_p\cong\overline{P_p}}.
\]
The condition $p\equiv1\pmod4$ ensures $-1$ is a square and adjacency is symmetric.

If $A$ is the adjacency matrix of $P_p$, define the permutation matrix by
$Pe_x=e_{\varphi(x)}$. The complement's adjacency matrix is $J-I-A$, and
the graph isomorphism gives
\[
PAP^{-1}=J-I-A.
\]
Similar matrices have the same characteristic polynomial, so $P_p$ and its complement have exactly the same eigenvalues with multiplicities. This is a second proof of the spectral symmetry without recomputing the Gauss sums.
\end{solution}

\begin{solution}{linear:ex:paleyramsey}
\idea Use the square residues modulo $17$. Translation and multiplication by non-zero squares preserve the graph, so we can move an arbitrary edge in a hypothetical four-clique to $\{0,1\}$. Then check the three common neighbours of these two vertices.

\writeup The set of non-zero squares modulo $17$ is
\[
S=\{1,2,4,8,9,13,15,16\}=\{\pm1,\pm2,\pm4,\pm8\}.
\]
The Paley graph $P_{17}$ joins two distinct residues when their difference is in $S$.

A translation $x\mapsto x+b$ preserves all differences. Scaling by a non-zero square also preserves squareness. Therefore, given any edge $\{x,y\}$, its difference $y-x$ is a square and the affine map
\[
z\longmapsto\frac{z-x}{y-x}
\]
is a graph automorphism sending that edge to $\{0,1\}$. If a four-clique existed, after such a transformation it would contain $0$ and $1$.

A third vertex $z$ must be adjacent to both $0$ and $1$, so both $z$ and $z-1$ are in $S$. Inspecting the eight residues above gives exactly
\[
z\in\{2,9,16\}.
\]
None of these three vertices is joined to another: their pairwise differences, up to sign, are $7,14,7$, and neither $7$ nor $14$ belongs to $S$. Thus we cannot add \emph{two} of them to $\{0,1\}$. No four-clique exists.

The three vertices $\{0,1,2\}$ do form a triangle because their non-zero differences are $\pm1$ or $\pm2$, all squares. Therefore
\[
\boxed{\omega(P_{17})=3}.
\]
By Exercise~\ref{linear:ex:paleyself}, the graph is isomorphic to its complement. Hence the complement also has clique number $3$, which means $P_{17}$ has no independent set of four vertices.

Colour the edges of the complete graph on $17$ vertices red for Paley edges and blue for non-edges. Neither colour contains a $K_4$. By the definition of Ramsey number,
\[
\boxed{R(4,4)\ge18}.
\]
Here $R(4,4)$ is the least number of vertices forcing a monochromatic
four-clique in every red/blue edge colouring. Such a finite number exists:
the same definition for $R(a,b)$, with $R(1,b)=R(a,1)=1$, satisfies, for $a,b\ge2$,
$R(a,b)\le R(a-1,b)+R(a,b-1)$. At a vertex of a complete graph on that
many vertices, either at least $R(a-1,b)$ incident edges are red or at least
$R(a,b-1)$ are blue; otherwise their total would be at most two less than
the vertex count. Apply the smaller Ramsey property in that neighbourhood
and either obtain the required clique there or adjoin the chosen vertex.
Induction proves finiteness. The displayed colouring therefore gives the
claimed lower bound.
\end{solution}

\begin{solution}{linear:ex:plotkinsharp}
\idea Plotkin's inequality comes from two estimates of the total distance among all ordered pairs of codewords. If the final bound is attained, both estimates must be equalities. The first forces every pairwise distance to be the minimum distance; the second forces every coordinate to be balanced.

\writeup An empty code cannot attain the positive bound. Let
$C\subseteq\{0,1\}^n$ be nonempty, with $M=|C|$ and minimum distance
at least $d$, where $d\ge1$ and $2d>n$. Let
\[
D=\sum_{x,y\in C}d_H(x,y)
\]
be the total Hamming distance over \emph{ordered} pairs. For $x\ne y$, the distance is at least $d$, while $d_H(x,x)=0$. Thus
\[
D\ge M(M-1)d.
\]
Equality holds only if \emph{every} distinct pair has distance exactly $d$.

For each coordinate $j$, let $w_j$ be the number of codewords with a $1$ in that coordinate. There are $w_j(M-w_j)$ unordered pairs disagreeing there, hence $2w_j(M-w_j)$ ordered pairs. Therefore
\[
D=2\sum_{j=1}^n w_j(M-w_j).
\]
The quadratic $w(M-w)$ is at most $M^2/4$, with equality only when $w=M/2$. Consequently
\[
D\le \frac{nM^2}{2}.
\]
Combining these bounds gives
\[
M(M-1)d\le \frac{nM^2}{2},
\]
and, dividing by $M>0$ and using $2d-n>0$,
\[
M\le\frac{2d}{2d-n}.
\]

Now suppose the code attains this last bound exactly. The first and last expressions in the chain are equal, so both intermediate inequalities must be equalities. Thus every non-zero pairwise distance equals $d$, and for every coordinate
\[
w_j(M-w_j)=M^2/4,
\]
forcing $w_j=M/2$. Hence each coordinate contains a $1$ in exactly half the
codewords. If $n\ge1$, this forces $M$ even. For $n=0$, the unique one-word
code attains the bound $1$, and both pairwise-distance and coordinate claims
are vacuous. These are necessary equality conditions; existence is a separate
question.
\end{solution}

\begin{solution}{linear:ex:lpthree}
\idea The basic Plotkin inequality only applies when $2d>n$, so it is silent at $n=8,d=4$. Delsarte's inequalities exploit more than the minimum distance: their first two Krawtchouk polynomials constrain the entire distance distribution. Combine those constraints with non-negative weights.

\writeup An empty code already has size at most $16$. Assume that
$C\subseteq\{0,1\}^8$ is nonempty, with minimum distance at least $4$,
and let $M=|C|$. Write $a_i$ for its average number of codewords at distance $i$ from a fixed codeword, averaged over all choices of that codeword. Then
\[
a_0=1,\qquad a_1=a_2=a_3=0,\qquad a_i\ge0.
\]
Also $M=1+\sum_{i=4}^8a_i$.

For length $8$, the first two Krawtchouk polynomials are
\[
K_1(i)=8-2i,\qquad
K_2(i)=\binom i2-i(8-i)+\binom{8-i}2.
\]
The Delsarte inequality $\sum_i a_iK_k(i)\ge0$ gives, at $k=1$,
\[
8-2a_5-4a_6-6a_7-8a_8\ge0,
\]
or
\[
2a_5+4a_6+6a_7+8a_8\le8. \tag{1}
\]
At $k=2$, the values for $i=0,4,5,6,7,8$ are $28,-4,-2,4,14,28$. Thus
\[
28-4a_4-2a_5+4a_6+14a_7+28a_8\ge0,
\]
which is
\[
4a_4+2a_5\le28+4a_6+14a_7+28a_8. \tag{2}
\]

Multiply (2) by $1/4$ and add it to (1). After moving all distance-distribution terms to the left, we obtain
\[
a_4+\frac52a_5+3a_6+\frac52a_7+a_8\le15.
\]
Every coefficient on the left is at least $1$, and every $a_i$ is non-negative. Therefore
\[
\sum_{i=4}^8a_i\le15,\qquad
\boxed{M\le16}.
\]

To attain the bound, label eight coordinates by the eight points of $\mathbb F_2^3$. Consider the truth tables of all affine functions
\[
f(x_1,x_2,x_3)=a_0+a_1x_1+a_2x_2+a_3x_3
\quad\text{with }a_i\in\mathbb F_2.
\]
There are $2^4=16$ coefficient choices, and their truth tables are distinct:
$f(0)=a_0$ and $f(e_i)+f(0)=a_i$ recover every coefficient.
The zero function gives the all-zero word; the constant-one function has weight $8$. Every non-constant affine function takes each value $0$ and $1$ exactly four times, because translation by a vector on which its linear part equals $1$ pairs zeros with ones. Thus every nonzero codeword has weight $4$ or $8$.

The set of affine truth tables is closed under binary addition. Hence the Hamming distance between two distinct words equals the weight of their non-zero difference and is at least $4$. This is the binary first-order Reed--Muller code $RM(1,3)$, a length-$8$ code of size $16$ and minimum distance $4$.

The construction and upper bound agree:
\[
\boxed{A(8,4)=16}.
\]
\end{solution}

\section*{Sources and further reading}

\begin{itemize}
\item T.~Hawkins, \emph{The origins of the theory of group characters}, Archive for History of Exact Sciences 7 (1971), 142--170; and
\emph{The Mathematics of Frobenius in Context}, Springer, 2013. The Dedekind--Frobenius correspondence of 1896.
\item P.~J.~Davis, \emph{Circulant Matrices}, Wiley, 1979; R.~M.~Gray, \emph{Toeplitz and circulant matrices: a review}, Foundations and
Trends in Communications and Information Theory 2 (2006), 155--239.
\item J.-P.~Serre, \emph{A Course in Arithmetic}, Springer, 1973, Chapter VI, Section 1 (characters of finite abelian groups); A.~Terras,
\emph{Fourier Analysis on Finite Groups and Applications}, Cambridge University Press, 1999.
\item R.~E.~A.~C.~Paley, \emph{On orthogonal matrices}, Journal of Mathematics and Physics 12 (1933), 311--320; R.~E.~Greenwood and
A.~M.~Gleason, \emph{Combinatorial relations and chromatic graphs}, Canadian Journal of Mathematics 7 (1955), 1--7; B.~Hanson and
G.~Petridis, \emph{Refined estimates concerning sumsets contained in the roots of unity}, Proceedings of the London Mathematical
Society 122 (2021), 353--358.
\item K.~Ireland and M.~Rosen, \emph{A Classical Introduction to Modern Number Theory}, 2nd ed., Springer, 1990, Chapter~6 (quadratic Gauss sums).
\item P.~Delsarte, \emph{An algebraic approach to the association schemes of coding theory}, Philips Research Reports Supplements 10
(1973); F.~J.~MacWilliams and N.~J.~A.~Sloane, \emph{The Theory of Error-Correcting Codes}, North-Holland, 1977;
R.~J.~McEliece, E.~R.~Rodemich, H.~Rumsey and L.~R.~Welch, \emph{New upper bounds on the rate of a code via the Delsarte--MacWilliams
inequalities}, IEEE Transactions on Information Theory 23 (1977), 157--166.
\item H.~Cohn and N.~Elkies, \emph{New upper bounds on sphere packings I}, Annals of Mathematics 157 (2003), 689--714; M.~S.~Viazovska,
\emph{The sphere packing problem in dimension 8}, Annals of Mathematics 185 (2017), 991--1015.
\item C.~Godsil and G.~Royle, \emph{Algebraic Graph Theory}, Springer, 2001, Chapter~10 (strongly regular graphs).
\end{itemize}
